\documentclass[11pt]{article}
\usepackage{fontspec}
\usepackage[french]{babel}
\usepackage{amsmath}
\usepackage{amssymb}
\usepackage{amsthm}
\usepackage{mathrsfs}
\usepackage{tikz-cd}
\usepackage{enumitem}

\newcommand{\printedpage}[1]{\marginpar{\footnotesize p.~#1}}
\newcommand{\msp}[1]{\marginpar{\footnotesize [#1]}}
% circled labels of the conditions, as printed
\newcommand{\circled}[1]{\tikz[baseline=(c.base)]\node[draw,circle,inner sep=1pt](c){#1};}
% "defined as" sign, as printed
\newcommand{\defeq}{\overset{\text{déf}}{=}}
% tilde printed after a closing parenthesis
\newcommand{\postilde}{\tilde{}}
\newcommand{\Fl}{\mathrm{Fl}}
\newcommand{\Cof}{\mathrm{Cof}}

\begin{document}

\printedpage{34}
% head: A. Grothendieck : Dérivateurs, Ch. XVII, Édition des Universités de Montpellier II et Paris VII
% foot: Version 9 février 2019
L'hypothèse \textit{stricte} pour une catégorie à cofibrations me paraît assez anodine, dans la plupart des cas à ma connaissance, elle est satisfaite. Mais même moyennant cette hypothèse, je ne vois aucun moyen de prouver l'existence de $u_!$ pour \textit{toute} flèche de Cat, et c'est une situation bien sûr très gênante. On voudrait absolument dégager des conditions qui assurent la validité de Dér 4 bis et de Dér 5 bis (existence de $u_!$ pour
\msp{page 71}
\textit{toute} flèche de Cat, et mode de calcul des fibres \dots).

\section*{4\quad Catégories à cofibrations spéciales}

\textbf{4.1.} Je suppose à présent que
\begin{itemize}[leftmargin=3em]
\item[\circled{a}] $\mathcal{M}$ est stable par $\varinjlim$ filtrantes, et
\item[\circled{b}] que ses objets sont accessibles,
\end{itemize}
de sorte qu'on est sous les conditions d'application du théorème de factorisation (XV, \S1). On suppose :
\begin{itemize}[leftmargin=3em]
\item[\circled{c}] donné $W, C \subset \Fl(\mathcal{M})$, satisfaisant les conditions $\Cof_1$, $\Cof_2$ de page 39. Pour $\Cof_3$, cf. plus bas.
\item[\circled{d}] $C$, $TC \defeq C \cap W$ \textit{stables par composition transfinie} (c'est le cas si la structure provient d'une structure de Quillen close $(W, C, F)$).
\end{itemize}
\[
TF \defeq C_* \;\subset\; \Fl(\mathcal{M})\ , \tag*{(4.1)}
\]
il résulte du théorème de factorisations que toute flèche $f$ de $\mathcal{M}$ se factorise en
\[
f \;=\; pi\ , \qquad p \in TF,\ i \in C\ ,
\]
\textit{pourvu} qu'on fasse l'hypothèse
\begin{itemize}[leftmargin=3em]
\item[\circled{e}] \mbox{}
\[
\exists\, C_0 \subset C,\ \textit{petite} \text{ partie de } C, \text{ telle que } C \subset \tilde{C}_0, \text{ i.e. } TF = (C_0)_*. \tag*{(4.3)}
\]
\end{itemize}
Sous ces conditions, si on suppose
\begin{itemize}[leftmargin=3em]
\item[\circled{f}] \mbox{}
\[
TF \subset W\ , \tag*{(4.4)}
\]
\end{itemize}
\msp{page 72}
alors l'axiome de factorisations $\Cof_3$ est satisfait. Inversement, si $\Cof_3$ est satisfait, et si $W$ est stable par facteurs directs, alors on a 4.4. Donc si $W$ est stable par facteurs directs, $\Cof_3$ équivaut à (4.4).

\textbf{4.2.} Il me semble que ces conditions, en pratique, ne seront guère jamais remplies, que si $W$ (qui est donné en tout premier lieu, satisfaisant $\Cof_0$ \textsuperscript{(}\footnote[18]{i.e. saturation faible.}\textsuperscript{)}), n'est pas stable par $\varinjlim$
\printedpage{35}
% head: A. Grothendieck : Dérivateurs, Ch. XVII, Édition des Universités de Montpellier II et Paris VII
% foot: Version 9 février 2019
filtrantes. À ce moment, pour construire la structure à cofibrations, on part de $C_0$ petite partie de $\Fl(\mathcal{M})$, prise telle que l'on ait
\[
(C_0)_* \;\subset\; W\ , \tag*{(4.5)}
\]
on prend pour $C$ au choix le saturé $C'_0$ de $C_0$ par cochangement de base et composition transfinie, ou $C = \tilde{C}_0 =$ ensemble des facteurs directs de $C'_0$. Les conditions $\Cof_1$ (stabilités de $C$), $\Cof_3$ (factorisation) sont satisfaites automatiquement, et il reste à s'occuper de $\Cof_2$, qui revient à la stabilité de $C{\cap}W$ par cochangement de base (la stabilité par composition et que ça contienne les isomorphismes, étant immédiates). C'est une condition délicate, et la seule façon standard que je vois pour l'assurer, c'est en supposant
\[
C_0 \;\subset\; \Cof_W\ ; \tag*{(4.6)}
\]
\msp{page 73}
la stabilité de $W$ par $\varinjlim$ filtrantes (donc aussi par facteurs directs) assure alors qu'on a également
\[
C \;\subset\; \Cof_W\ , \tag*{(4.7)}
\]
et par suite
\[
\left\{
\begin{aligned}
& C \cap W \subset \Cof_W \cap W = W^{\mathrm{univ}}, \\
& \text{d'où }\ TC \defeq C \cap W = C \cap W^{\mathrm{univ}}\ ,
\end{aligned}
\right. \tag*{(4.8)}
\]
ce qui assure la stabilité de $C \cap W$ par cochangement de base. La stabilité de $TC$ par composition transfinie résulte de celle de $C$, et de la stabilité de $W$ par composition transfinie.

\textbf{4.3.} Je reviens aux conditions générales\circled{a} à\circled{f} de 4.1, sans supposer nécessairement $W$ stable par $\varinjlim$ filtrantes. Je veux à présent étendre la structure à cofibrations aux $\mathcal{M}(I)$, $I$ quelconque dans Cat. Nous n'avons pas pour ceci (je crois) à utiliser les développements des paragraphes 2, 3, mais nous ferons le lien avec ceux-ci par la suite.

Désignons par $I_0$ la catégorie discrète ayant mêmes objets que $I$, d'où un foncteur canonique
\[
v = v_I \ :\ I_0 \longrightarrow I\ . \tag*{(4.9)}
\]
L'idée est de définir la structure de $\mathcal{M}(I)$ via la structure évidente de $\mathcal{M}(I_0) \simeq \mathcal{M}^{I_0}$, en jouant sur le couple de foncteurs adjoints
\[
\begin{tikzcd}
\mathcal{M}(I_0) \arrow[r, shift left, "\varphi = v_!"] & \mathcal{M}(I) \arrow[l, shift left, "\psi = v^*"]
\end{tikzcd}
\ .
\]
\msp{page 74}
Nous allons présenter les choses de façon un peu plus générale.

\textbf{Définition 4.4.} J'appelle structure de cofibrations \textit{spéciale} une structure $(\mathcal{M}, W, C)$ satisfaisant les conditions a) à f) de 4.1, que je reformule :
\begin{itemize}[leftmargin=3em]
\item[\circled{a}] $\mathcal{M}$ stable par petites limites inductives, et les objets de $\mathcal{M}$ sont accessibles.
\printedpage{36}
% head: A. Grothendieck : Dérivateurs, Ch. XVII, Édition des Universités de Montpellier II et Paris VII
% foot: Version 9 février 2019
\item[\circled{b}] $W$ est faiblement saturé (i.e. $\Cof_0$ de p. 39), $C$ stable par cochangement de base et par composition, $TC \defeq C{\cap}W$ stable par cochangement de base (et par composition automatiquement).
\item[\circled{b$'$}] $C$ et $TC$ sont stables même par composition \textit{transfinie.}
\item[\circled{c}] $\exists$ petite partie $C_0 \subset C$ telle que $C \subset \tilde{C}_0$ ($\Longrightarrow \tilde{C} = \tilde{C}_0$).
\item[\circled{d}] $TF \defeq C_* \subset W$ (ce qui implique, joint à\circled{a} à\circled{c}, l'axiome de factorisation $\Cof_3$ de page 39).
\end{itemize}

\textbf{4.5.} Nous allons donner une recette générale d'extension d'une telle structure, quand on dispose d'un couple de foncteurs adjoints
\[
\begin{tikzcd}
\mathcal{M} \arrow[r, shift left, "v_!"] & \mathcal{M}' \arrow[l, shift left, "v^*"]
\end{tikzcd}
\ ,
\]
$\mathcal{M}$ étant muni d'une structure de cofibrations spéciale $(W, C)$, et $\mathcal{M}'$ satisfaisant la condition liminaire\circled{a}. On définira alors les parties $W'$, $C'$ de $\Fl(\mathcal{M}')$
\msp{page 75}
par
\[
\left\{
\begin{aligned}
W' \;&=\; (v^*)^{-1}(W) \\
C' \;&=\; \text{clôture stable par cochangement de base} \\
&\phantom{\;=\;}\ \ \text{et composition transfinie de } v_!(C),
\end{aligned}
\right. \tag*{(4.5.1)}
\]
donc
\[
v_!(C) \;\subset\; C' \;\subset\; (v_!(C))\postilde \;=\; \tilde{C}'\ . \tag*{(4.5.2)}
\]
Je dis que $(W', C')$, et au choix $(W', \tilde{C}')$, est une structure de cofibrations spéciale sur $\mathcal{M}'$. On va vérifier les différentes conditions de 4.4. La condition a) est déjà acquise.

b) \textsuperscript{(}\footnote[19]{\textbf{N.B.} Il vaut mieux commencer par prouver c).}\textsuperscript{)}. La saturation faible de $W'$ résulte de celle de $W$, via (4.5.1). La stabilité de $C'$ par cochangement de base et par composition (même transfinie) est acquise par définition de $C'$, et itou pour $\tilde{C}'$ (étant Q-close à gauche). Stabilité de $TC' \defeq C'{\cap}W'$ par cochangement de base ? Il suffit de voir que si on a un carré cocartésien
\[
\begin{tikzcd}
X'_0 \arrow[r, "i'"] \arrow[d] \arrow[dr, phantom, "\text{\scriptsize cocart.}"] & X'_1 \arrow[d] \\
X'_2 \arrow[r, "j'"] & X'_3
\end{tikzcd}
\]
dans $\mathcal{M}'$, avec $i'$ dans $TC' = C' \cap W'$, alors $j' \in W'$ (car on sait déjà $j' \in C'$), i.e.
\marginpar{$||$}% marginal sign || printed in the left margin of this line
$v^*(j') \in W'$. Nous allons supposer que $v^*$ \textit{commute aux sommes amalgamées,} et on est donc ramené à prouver que $v^*(i') \in TC$. Donc % [sic] $v^*(j') \in W'$ -> $v^*(j') \in W$, since $W' = (v^*)^{-1}(W)$
\printedpage{37}
% head: A. Grothendieck : Dérivateurs, Ch. XVII, Édition des Universités de Montpellier II et Paris VII
% foot: Version 9 février 2019
\msp{page 76}
on doit prouver le

\textbf{Lemme 4.5.2.} \textit{Sous les conditions} (4.5.3) \textit{et} (4.5.4) \textit{ci-dessous, on a}
\[
v^*(C') \;\subset\; C\ , \qquad v^*(TC') \;\subset\; TC\ ,
\]
\textit{et si $C = \tilde{C}$, on a $v^*(\tilde{C}') \subset C$, $v^*(T\tilde{C}') \subset TC$.}

Pour $v^*(C') \subset C$, il faut d'abord prouver
\[
v^*(v_!(C)) \;\subset\; C\ , \tag*{(4.5.3)}
\]
et je ne vois aucun moyen -- il faut le poser comme hypothèse. De plus, que $v^*$ commute aux $\varinjlim$ quelconques (cochangement de base déjà admis -- à présent on a besoin des limites inductives filtrantes)
\[
v^* \text{ commute aux petites } \varinjlim \text{ quelconques} \tag*{(4.5.4)}
\]
pour conclure de (4.5.3) que $v^*(C' =$ clôture de $v_!(C)$ pour cochangement de base et composition transfinie) est également dans $C$. Passant à $v^*(TC') \subset TC$, par ce qui précède, $v^*(C') \subset C$, il suffit de prouver $v^*(TC') \subset W$, mais $v^*(TC') \subset v^*(W') \subset W$ par définition de $W'$.

Supposons $C = \tilde{C}$, i.e. $C$ stable par facteurs directs, et prouvons $v^*(\tilde{C}') \subset C$, $v^*(T\tilde{C}') \subset TC$. Par c) plus bas, $\tilde{C}'$ est formé des facteurs directs
\msp{page 77}
d'éléments de $C'$, donc $v^*(C') \subset C$ implique $v^*(\tilde{C}') \subset C$ puisque $C$ stable par facteurs directs.

Prouvons enfin $v^*(T\tilde{C}') \subset TC \defeq C \cap W$. Comme $T\tilde{C}' = \tilde{C}' \cap W'$, cela résulte de $v^*(\tilde{C}') \subset C$, et de $v^*(W') \subset W$.

Prouvons b$'$). La stabilité de $C'$ et de $\tilde{C}'$ par composition transfinie est acquise, il reste à prouver la stabilité de $TC' = C' \cap W'$ et de $T\tilde{C}' = \tilde{C}' \cap W'$. Il suffit de voir que le composé transfinie $f$ en question est dans $W'$, donc que $v^*(f) \in W$, mais cela résulte du fait que $v^*$ commute à la composition transfinie par (4.5.4), et que $v^*$ transforme $TC'$ et $T\tilde{C}'$ dans $TC$.

\textit{Prouvons} c) pour $C'$ (donc aussi pour $\tilde{C}'$). Soit $C_0$ une petite partie de $C$ telle que $C \subset \tilde{C}_0$, i.e. $C$ est contenu dans la clôture de $C_0$ par les trois opérations de cochangement de base, de composition transfinie, des facteurs directs. On voit alors que $v_!(C)$ est dans la clôture de $v_!(C_0)$ pour les opérations, donc aussi $C'$, et on peut alors prendre $C'_0 = v_!(C_0)$.

\textit{Prouvons enfin} d), c'est la même assertion pour $C'$ et pour $\tilde{C}'$, puisque
\[
C'_* \;=\; \tilde{C}'_*\ (\ \defeq\ TF'\ )\ ,
\]
\msp{page 78}
il faut donc prouver
\[
TF' \defeq C'_* \;\subset\; W'\ .
\]
Cela va résulter du
\printedpage{38}
% head: A. Grothendieck : Dérivateurs, Ch. XVII, Édition des Universités de Montpellier II et Paris VII
% foot: Version 9 février 2019

\textbf{Lemme 4.5.3.} \textit{On a $TF' = (v^*)^{-1}(TF)$.}

C'est formel par définition de $TF'$ comme
\[
TF' \;=\; (v_!(C))_*\ ,
\]
de $TF$ comme
\[
TF \;=\; C_*\ ,
\]
et de l'adjonction de $v_!$, $v^*$. Comme $TF \subset W$, on en conclut
\[
TF' \;=\; (v^*)^{-1}(TF) \;\subset\; (v^*)^{-1}(W) = W'\ ,
\]
q.e.d.

Donc on a prouvé ceci :

\textbf{Proposition 4.6.} \textit{Considérons un couple de foncteurs adjoints}
\[
\begin{tikzcd}
\mathcal{M} \arrow[r, shift left, "v_!"] & \mathcal{M}' \arrow[l, shift left, "v^*"]
\end{tikzcd}
\ ,
\]
\textit{où $\mathcal{M}$ est muni d'une structure de cofibrations} spéciale $(W, C)$\textit{. On suppose de plus les conditions suivantes satisfaites.}
\[
\begin{aligned}
&1^\circ)\quad \mathcal{M}' \textit{ stable par petites } \varinjlim\ \textit{, ses éléments sont accessibles.} \\
&2^\circ)\quad v^* \textit{ commute aux petites } \varinjlim\ . \\
&3^\circ)\quad v^*(v_!(C)) \subset C.
\end{aligned} \tag*{(4.6.1)}
\]
\msp{page 79}
\textit{Considérons les parties $W'$, $C'$ de} $\Fl(\mathcal{M}')$ \textit{définies par}
\[
\left\{
\begin{aligned}
W' \;&=\; (v^*)^{-1}(W), \\
C' \;&=\; \textit{clôture de } v_!(C) \textit{ par cochangement de base et par composition} \\
&\phantom{\;=\;}\ \textit{transfinie.}
\end{aligned}
\right. \tag*{(4.6.2)}
\]
\textit{Alors $(\mathcal{M}', W', C')$ est une structure à cofibrations spéciale. Si de plus $C = \tilde{C}$, i.e. $C$ est stable par facteurs directs, alors $(\mathcal{M}', W', \tilde{C}')$ est également une structure à cofibrations spéciale} \textsuperscript{(}\footnote[20]{\textbf{N.B.} C'est là un fait général pour les catégories à cofibrations spéciales, et devrait être explicité plus bas (cf. 5.1). L'hypothèse sur $C$ est inutile ici (ou plutôt, il suffit au lieu de cela que $W$ soit stable par facteurs directs \textit{et} par composition transfinie, ce qui n'est pas une hypothèse plus réjouissante \dots).}\textsuperscript{)}.

\textbf{4.7.} J'aimerais, en vue d'énoncés du type 3.7 (établissant l'existence de foncteurs du type $u_!^W$), avoir une prise sur $TC' = C'{\cap}W'$, et sur $T\tilde{C}' = \tilde{C}'{\cap}W'$, de façon précise, j'aimerais avoir la relation
\[
v_!(TC)\ \subset TC' \subset v_!(TC)\postilde\ , \tag*{(4.7.1)}
\]
voire la relation légèrement plus forte
\[
\underbrace{v_!(TC)\ \subset TC'}_{} \ \overset{\text{tautol.}}{\subset}\ \underbrace{T\tilde{C}' \;\subset\; (v_!(TC))\postilde}_{}\ , \tag*{(4.7.2)}
\]

\end{document}
